% Source: https://artofproblemsolving.com/community/c3h1403737_something_i_was_curious_about__
% Stable identity: aops-topic-1403737-q01
Recall that $\dfrac{a}{b}$ is defined (for $b\ne0$) as the unique number $q$ with $bq=a$, so $b\cdot\dfrac{a}{b}=a$.

\textbf{If $\dfrac{a}{b}=\dfrac{c}{d}$}, multiply both sides by $bd$:
\[
bd\cdot\dfrac{a}{b}=d\left(b\cdot\dfrac {a}{b}\right)=da,\qquad bd\cdot\dfrac{c}{d}=b\left(d\cdot\dfrac {c}{d}\right)=bc,
\]
so $ad=bc$.

\textbf{If $ad=bc$}, divide both sides by the nonzero number $bd$ (that is, multiply by $\dfrac{1}{bd}$):
\[
\dfrac{ad}{bd}=\dfrac{a}{b},\qquad \dfrac{bc}{bd}=\dfrac{c}{d},
\]
so $\dfrac {a}{b}=\dfrac {c}{d}$.

The conditions are needed because division by zero is undefined; moreover, without them the converse fails: $a=1$, $b=0$, $c=2$, $d=0$ satisfy $ad=bc=0$ although $\dfrac{1}{0}$ and $\dfrac{2}{0}$ mean nothing. So cross-multiplication is not an axiom; it follows from the definition of division.
